\section{Pessoa 2 -- Solicitacao de Licenca/Folga}

\subsection{Introducao}
Esta secao descreve a contribuicao individual da Pessoa 2 no Trabalho T2 de Verificacao e Validacao de Software. O sistema sob teste e o Orange Human Resource Management (OrangeHRM), um sistema de gestao de recursos humanos que inclui cadastro de funcionarios, usuarios, licencas, saldos de licenca e fluxos administrativos de aprovacao.

O codigo-fonte oficial esta disponivel em \url{https://github.com/orangehrm/orangehrm}. Para analise local, foi clonado o commit \texttt{56e23b3b09e7af29317a0943523b825843fff527}, referente a mensagem \texttt{OHRM5X-2669: Update Dockerfile for OrangeHRM 5.9 (\#1962)}. A copia local foi armazenada em \texttt{external/orangehrm/} e ignorada pelo Git do trabalho. O servidor demonstrativo usado no E2E foi \url{https://opensource-demo.orangehrmlive.com}.

\subsection{Objetivo dos Testes}
O objetivo foi verificar a jornada em que um funcionario solicita um periodo de licenca/folga e um administrador aprova ou rejeita a solicitacao. A verificacao envolveu tres frentes:

\begin{itemize}
    \item analise dos testes ja existentes no fonte oficial do OrangeHRM;
    \item testes unitarios e de integracao em uma modelagem auxiliar no repositorio da turma;
    \item testes E2E reais com Playwright contra o servidor demo do OrangeHRM.
\end{itemize}

\subsection{Descricao da Jornada de Usuario}
A jornada de usuario selecionada foi:

\begin{enumerate}
    \item o administrador cria ou disponibiliza um funcionario com usuario de login;
    \item o administrador concede saldo de licenca ao funcionario;
    \item o funcionario faz login no OrangeHRM;
    \item o funcionario acessa o modulo \texttt{Leave} e solicita licenca;
    \item o sistema registra a solicitacao como pendente;
    \item o administrador faz login;
    \item o administrador acessa \texttt{Leave List};
    \item o administrador aprova ou rejeita a solicitacao;
    \item a solicitacao deixa de aparecer na lista de pendentes.
\end{enumerate}

\subsection{Analise do Fonte Oficial}
No fonte oficial, a funcionalidade esta localizada principalmente em \texttt{src/plugins/orangehrmLeavePlugin}. Os principais artefatos identificados foram:

\begin{itemize}
    \item \texttt{Api/MyLeaveRequestAPI.php}: API usada pelo funcionario para criar e consultar as proprias solicitacoes;
    \item \texttt{Api/EmployeeLeaveRequestAPI.php}: API usada por administrador ou supervisor para consultar e alterar solicitacoes de funcionarios;
    \item \texttt{Service/LeaveRequestService.php}: regras de servico relacionadas a solicitacoes;
    \item \texttt{Dao/LeaveRequestDao.php}: persistencia e consulta das solicitacoes;
    \item \texttt{entity/LeaveRequest.php} e \texttt{entity/Leave.php}: entidades de dominio;
    \item \texttt{test/Api/MyLeaveRequestAPITest.php}, \texttt{test/Api/EmployeeLeaveRequestAPITest.php}, \texttt{test/Service/LeaveRequestServiceTest.php}, \texttt{test/Dao/LeaveRequestDaoTest.php} e \texttt{test/Api/LeaveOverlapAPITest.php}: testes existentes no projeto oficial.
\end{itemize}

A cobertura oficial observada e forte em API, servico, DAO e uso de fixtures. Como ponto de melhoria, para o escopo deste trabalho, foi acrescentado um teste E2E de navegador cobrindo a jornada completa de usuario.

\subsection{Escopo da Contribuicao Individual}
A contribuicao individual criou:

\begin{itemize}
    \item configuracao Playwright em \texttt{playwright.config.js};
    \item teste E2E em \texttt{tests/e2e/orangehrm-leave-request.spec.js};
    \item documentacao em \texttt{PESSOA2\_LICENCA\_TESTES.md} e neste arquivo \texttt{.tex};
    \item modelagem Java auxiliar com \texttt{LeaveRequest}, \texttt{LeaveRequestService}, \texttt{LeaveType} e \texttt{LeaveRequestStatus};
    \item testes JUnit auxiliares unitarios, de integracao e de sistema para as regras da jornada.
\end{itemize}

A modelagem Java nao substitui o OrangeHRM oficial; ela serve apenas como artefato auxiliar para demonstrar tecnicas de teste unitario e de integracao dentro do repositorio da turma.

\subsection{Tecnicas de Teste Usadas}
Foram utilizadas as seguintes tecnicas:

\begin{itemize}
    \item particionamento de equivalencia para entradas validas e invalidas;
    \item analise de valor limite para datas e duracao de periodo;
    \item teste de transicao de estados para pendente, aprovado e rejeitado;
    \item teste negativo para solicitacao inexistente, decisao sem administrador e periodo sobreposto;
    \item teste de integracao entre servico e entidade auxiliar;
    \item teste E2E de fluxo real com navegador.
\end{itemize}

\subsection{Casos de Teste}
\begin{table}[h]
\centering
\begin{tabular}{|p{2.1cm}|p{6.0cm}|p{3.0cm}|p{3.0cm}|}
\hline
\textbf{Caso} & \textbf{Cenario} & \textbf{Resultado Esperado} & \textbf{Nivel} \\
\hline
CT-LIC-01 & Solicitacao valida na modelagem auxiliar & Status pendente & Unitario \\
\hline
CT-LIC-02 & Funcionario vazio, tipo nulo ou datas nulas & Excecao de validacao & Unitario \\
\hline
CT-LIC-03 & Data final anterior a inicial & Excecao de periodo invalido & Unitario \\
\hline
CT-LIC-04 & Periodo de um dia & Duracao igual a 1 & Unitario \\
\hline
CT-LIC-05 & Administrador aprova solicitacao pendente & Status aprovado & Integracao/Sistema \\
\hline
CT-LIC-06 & Administrador rejeita solicitacao pendente & Status rejeitado & Integracao/Sistema \\
\hline
CT-LIC-07 & Periodo sobreposto ativo & Nova solicitacao bloqueada & Integracao \\
\hline
CT-LIC-08 & Solicitacao rejeitada seguida de nova solicitacao & Nova solicitacao permitida & Integracao \\
\hline
CT-LIC-09 & Funcionario real solicita licenca no demo e admin aprova & Solicitacao sai da lista de pendentes & E2E \\
\hline
CT-LIC-10 & Funcionario real solicita licenca no demo e admin rejeita & Solicitacao sai da lista de pendentes & E2E \\
\hline
\end{tabular}
\caption{Casos de teste da jornada de solicitacao de licenca/folga.}
\end{table}

\subsection{Testes Unitarios}
Os testes unitarios foram implementados em \texttt{LeaveRequestTest}. Eles validam campos obrigatorios, tipo de licenca, datas, duracao, aprovacao, rejeicao e impedimento de nova decisao sobre solicitacao ja finalizada.

\subsection{Testes de Integracao}
Os testes de integracao foram implementados em \texttt{LeaveRequestServiceIntegrationTest}. Eles exercitam criacao, consulta, listagem por funcionario, aprovacao, rejeicao, bloqueio de sobreposicao ativa, liberacao apos rejeicao e solicitacao inexistente.

\subsection{Testes de Sistema/E2E}
O teste E2E real foi implementado em \texttt{tests/e2e/orangehrm-leave-request.spec.js}, usando Playwright. Ele roda contra o servidor demo do OrangeHRM e executa dois cenarios:

\begin{itemize}
    \item funcionario solicita licenca e administrador aprova;
    \item funcionario solicita licenca e administrador rejeita.
\end{itemize}

O E2E cria um funcionario temporario, cria usuario de login para esse funcionario, adiciona entitlement de licenca, autentica como funcionario, solicita licenca e autentica novamente como administrador para decidir a solicitacao. O teste fica protegido por \texttt{ORANGEHRM\_E2E\_ENABLED=true}, pois altera dados do servidor demo.

\subsection{Como Executar}
Para executar os testes Java:

\begin{verbatim}
cd tf-verival
mvn test
\end{verbatim}

Para executar o E2E real:

\begin{verbatim}
npm install
ORANGEHRM_E2E_ENABLED=true npx playwright test
\end{verbatim}

\subsection{Resultados Obtidos}
O baseline inicial do Maven executou 12 testes da parte existente de funcionarios, com 0 falhas, 0 erros e 0 ignorados.

A execucao Java final executou 35 testes, com 0 falhas, 0 erros e 1 ignorado. O teste ignorado foi o preflight JUnit do E2E remoto, mantido desabilitado por padrao.

A execucao Playwright contra o OrangeHRM demo executou 2 testes E2E reais, ambos aprovados:

\begin{itemize}
    \item funcionario solicita licenca e admin aprova a solicitacao;
    \item funcionario solicita licenca e admin rejeita a solicitacao.
\end{itemize}

Tempo observado da execucao E2E: aproximadamente 1,6 minuto.

\subsection{Defeitos Encontrados}
Nao foram encontrados defeitos nos cenarios executados. Os dois caminhos E2E passaram no servidor demo no momento da execucao. Como o ambiente demonstrativo e compartilhado e mutavel, esse resultado nao garante ausencia de defeitos em execucoes futuras.

\subsection{Limitacoes}
Os testes PHP oficiais do OrangeHRM foram analisados, mas nao executados neste repositorio, pois dependem do ambiente completo do OrangeHRM. O E2E usa o servidor demo publico, que pode sofrer alteracoes de dados, credenciais e desempenho. Alem disso, os testes E2E criam dados temporarios no demo.

\subsection{Registro do Uso de IA}
Este material foi elaborado com auxilio da ferramenta ChatGPT/Codex e deve ser revisado pelo aluno antes da entrega. A responsabilidade pela validacao final, registro de Issue, commits e adequacao ao formato da disciplina permanece com o autor da contribuicao.

\subsection{Referencias}
\begin{itemize}
    \item OrangeHRM. Codigo-fonte oficial. Disponivel em: \url{https://github.com/orangehrm/orangehrm}.
    \item OrangeHRM. Servidor demonstrativo. Disponivel em: \url{https://opensource-demo.orangehrmlive.com}.
    \item OrangeHRM. Site institucional. Disponivel em: \url{https://www.orangehrm.com}.
    \item Playwright. Documentacao. Disponivel em: \url{https://playwright.dev}.
    \item JUnit 5. Documentacao. Disponivel em: \url{https://junit.org/junit5/docs/current/user-guide/}.
\end{itemize}
